The local transition table (1.771) and logistic regression (1.723) are close to the brief's test references of 1.81 and 1.75 \cite{brief}. This is a sanity check that the validation cutoffs capture a comparable level of difficulty.

With a cap of 600 rounds, the patience rule triggered in 5 of 5 tuning trials. The inner holdout selected leaves7\_child220\_l2\_10 at its best iteration of 254, after evaluating 274 rounds. Adding back the 196-category event pair changed validation loss by +0.0219. With the pair excluded in both models, dropping the six individual lags changed loss by +0.0080. Dropping sensor history changed loss by +0.0024, with an interval spanning zero. The remaining features include current sensor readings. The sensor history did not show a clear incremental effect. Using four cutoffs per fitting machine changed loss by +0.0345 at the same tree settings.

Temperature scaling fitted \(T=0.9955\) and changed loss by -0.0001, with a paired interval excluding zero. Temperature scaling gave a statistically detectable but negligible gain (about 0.00011 log-loss points). The calibration split selected temperature scaling. For the three rare target classes, the share assigned under 0.5\% probability for the true event was 11.6\% raw and 11.9\% after temperature scaling. Prior blending fitted \(\alpha=0\), so it reproduced the raw probabilities and is omitted from the table. The paired intervals describe these fitted models and do not adjust for multiple comparisons or include variation across refits.
